\documentclass[10pt,a4paper]{article}
\usepackage{amsmath,amssymb,amsthm,graphicx}
\usepackage{amsmath,amssymb,amsthm}
\usepackage{hyperref}
\newtheorem{theorem}{Theorem}
\newtheorem{lemma}[theorem]{Lemma}
\newtheorem{claim}[theorem]{Claim}
\newtheorem{definition}[theorem]{Definition}
\newcommand{\F}{\mathbb F}
\newcommand{\N}{\mathbb N}
\newcommand{\vx}{\mathbf x}
\DeclareMathOperator{\fdeg}{fdeg}
\DeclareMathOperator{\Char}{char}
\DeclareMathOperator{\depth}{\mathsf{depth}}
\newcommand{\set}[1]{\left\{#1\right\}}
\newcommand{\abs}[1]{\left|#1\right|}
\newcommand{\floor}[1]{\left\lfloor#1\right\rfloor}
\newcommand{\inparen}[1]{\left(#1\right)}
\providecommand{\subsectionautorefname}{Section}
\renewcommand{\subsectionautorefname}{Section}
\title{A quadratic vertex lower bound for algebraic branching programs via depth reduction robust to structured errors}
\author{Prerona Chatterjee, Mrinal Kumar, Adrian She, Ben Lee Volk}
\date{doi:10.4230/LIPIcs.CCC.2020.2; CC BY 3.0 Unported}
\begin{document}
\maketitle
\noindent\textit{Editorial note.} Editable transcription of the formally published CC BY3.0 Unported quadratic ABP result. Nonexclusive nativev1/v2 used privately only; public files contain only the selected formal model, Sections2–3 and actual cited bibliography. Formal publication additions/wording are retained, including the formal-degree monotonicity sentence, Lemma4 degree range, Section2.2 dimension properties, variety-inclusion argument and explicit d0Delta formal-degree bound. No subsequent formula proof is included or relicensed. Original author equations and constants retained without mathematical correction. Minimal macro presentation and original global theorem numbering restored; no input/include, unsupported fonts or unused native packages. All production/validation only on Yupeng; root owns compilation.\par\medskip
\noindent\textit{Allowed standard prerequisites (editorial).}\begin{itemize}
\item Affine varieties and dimension bounds; classical homogeneous ABP lower bound and finite derivative-zero variety lemma. FormalSection2.2 standard dimension properties and cited Kumar claim retained; source-local decomposition, robustness proof, shrinking and iterative reduction included completely
\end{itemize}\medskip
\section*{Model and main theorem (original Section 1)}
\begin{definition}[Algebraic Branching Programs]\label{def:abp}
An Algebraic Branching Program (ABP) is a layered graph where each edge is labeled by an affine linear form and the first and the last layer have one vertex each, called the ``start'' and the ``end'' vertex respectively.
The polynomial computed by an ABP is equal to the sum of the weights of all paths from the start vertex to the end vertex in the ABP, where the weight of a path is equal to the product of the labels of all the edges on it.
The size of an ABP is the number of vertices in it.
\end{definition}
A more general definition, which we sometimes adopt, is to allow every edge to be labeled by an arbitrary polynomial of degree at most $\Delta$. In this case we refer to the model as an ABP with edge labels of degree at most $\Delta$. Thus, the common case is $\Delta=1$, but our results are meaningful even when $\Delta=\omega(1)$.

\begin{theorem}\label{thm:intro main}
Let $\F$ be a field and $n\in\N$ such that $\Char(\F)\nmid n$. Then any algebraic branching program over $\F$ computing the polynomial $\sum_{i=1}^n x_i^n$ is of size at least $\Omega(n^2)$.
When the ABP's edge labels are allowed to be polynomials of degree at most $\Delta$, our lower bound is $\Omega(n^2/\Delta)$.
\end{theorem}

\setcounter{section}{1}
\setcounter{theorem}{3}
\section{Notations and Preliminaries}
All logarithms in the paper are base 2.

We use some standard graph theory terminology: If $G$ is a directed graph and $(u,v)$ is an edge, $v$ is called the \emph{head} of the edge and $u$ the \emph{tail}. Our directed graphs are always acyclic with designated source vertex $s$ and sink vertex $t$. The \emph{depth} of a vertex $v$, denoted $\depth(v)$, is the length (in edges) of a longest path from $s$ to $v$. The depth of the graph, denoted by $\depth(G)$, is the depth of $t$.

For any two vertices $u$ and $v$ in an ABP, the polynomial computed between $u$ and $v$ is the sum of weights of all paths between $u$ and $v$  in the ABP. We denote this by $[u,v]$. 

The formal degree of a vertex $u$ in an ABP denoted $\fdeg(u)$, is defined inductively as follows: If $s$ is the start vertex of the ABP, $\fdeg(s)=0$. If $u$ is a vertex with incoming edges from $u_1, \ldots, u_k$, labeled by non-zero polynomials $\ell_1, \ldots, \ell_k$, respectively, then
\[
\fdeg(u) = \max_{i \in [k]} \set{ \deg(\ell_i) + \fdeg(u_i) }.
\]
It follows by induction that for every vertex $u$, $\deg ([s,u]) \le \fdeg(u)$ (however, cancellations can allow for arbitrary gaps between the two). Also, note that the formal degree of vertices is monotonically non-decreasing along any path from the source vertex to the sink vertex. The formal degree of the ABP is the maximal formal degree of any vertex in it.

We sometimes denote by $\vx$ the vector of variables $(x_1, \ldots, x_n)$, where $n$ is understood from the context. Similarly we use $\mathbf{0}$ to denote the $n$-dimensional vector $(0,0,\ldots,0)$.

	

\subsection{A decomposition lemma}
The following lemma gives a decomposition of a (possibly unlayered) ABP in terms of the intermediate polynomials it computes. Its proof closely resembles that of Lemma 3.5 of \cite{K19}. For completeness we prove it here for a slightly more general model.

\begin{lemma}[\cite{K19}]\label{lem:structure}
	Let $\mathcal{B}$ be a (possibly unlayered) algebraic branching program which computes a degree $d$ polynomial $P\in\F[x_1,\ldots,x_n]$, and has formal degree $d$. Further, assume that the edges of $\mathcal{B}$ are labelled by arbitrary polynomials of degree at most $\Delta$, where $1\le\Delta\le d/2$. Set $d'=\floor{d/\Delta}$.
		
	 For any $i \in \set{1, 2, \ldots ,d'-1}$, let $S_i = \set{u_{i,1}, u_{i,2}, \ldots, u_{i,m}}$ be the set of all vertices in $\mathcal{B}$ whose formal degree is in the interval $[i \Delta, (i+1)\Delta)$.
	 
	 Then, there exist polynomials $Q_{i,1}, Q_{i,2}, \ldots, Q_{i,m}$ and $R_i$, each of degree at most $d-1$ such that
		\[P = \sum_{j=1}^{m} [s, u_{i,j}] \cdot Q_{i,j} + R_i\, .\]
\end{lemma}

\begin{proof}
Fix $i$ as above and set $S_i = \set{u_{i,1}, u_{i,2}, \ldots, u_{i,m}}$ as above (observe that since each edge label is of degree at most $\Delta$, $S_i$ is non empty). Further suppose, without loss of generality, that the elements of $S_i$ are ordered such that there is no directed path from $u_{i,j}$ to $u_{i,j'}$ for $j' > j$.

Consider the unlayered ABP $\mathcal{B}_1$ obtained from $\mathcal{B}$ by erasing all incoming edges to $u_{i,1}$, and multiplying all the labels of the outgoing edges from $u_{i,1}$ by a new variable $y_1$. The ABP $\mathcal{B}_1$ now computes a polynomial of the form
\[
	P'(y_1, x_1, \ldots, x_n) = y_1 \cdot Q_{i,1} + R_{i,1}
\]
where $P = P'([s,u_{i,1}], x_1, \ldots, x_n)$. $R_{i,1}$ is the polynomial obtained from $\mathcal{B}_1$ by setting $y_1$ to zero, or equivalently, removing $u_{i,1}$ and all its outgoing edges. We continue in the same manner with $u_{i,2}, \ldots, u_{i,m}$ to obtain
\[
P = \sum_{j=1}^m [s,u_{i,j}] \cdot Q_{i,j} + R_i.
\]
Indeed, observe that since there is no path from $u_{i,j}$ to $u_{i,j'}$ for $j' > j$, removing $u_{i,j}$ does not change $[s,u_{i,j'}]$. The bound on the degrees of $Q_{i,j}$ is immediate from the fact that the formal degree of the ABP is at most $d$ and $\fdeg(u_{i,j}) \ge 1$. It remains to argue the $\deg(R_i) \le d-1$.

The polynomial $R_i$ is obtained from $\mathcal{B}$ by erasing all the vertices in $S_i$ and the edges touching them. We will show that every path in the corresponding ABP computes a polynomial of degree at most $d-1$. Let $s=v_1, v_2, \ldots, v_r=t$  be such a path, which is also a path in $\mathcal{B}$. Let $v_k$ be the minimal vertex in the path whose degree (in $\mathcal{B}$) is at least $(i+1)\Delta$
(if no such $v_k$ exists, the proposition follows).
As $v_{k-1} \not \in S_i$, the formal degree of $v_{k-1}$ is at most $i \Delta - 1$.
The degree of the polynomial computed by this path is thus at most $i\Delta - 1 + \Delta + D = (i+1)\Delta - 1 + D$, where $D$ is the degree of product of the labels on the path $v_k, v_{k+1}, \ldots, t$. To complete the proof, it remains to be shown that $D \le d - (i+1)\Delta$.

Indeed, if $D \ge d - (i+1)\Delta + 1$ then since the degree of $v_k$ is at least $(i+1)\Delta$, there would be in $\mathcal{B}$ a path of formal degree at least $(i+1)\Delta +D \ge d + 1$, contradicting the assumption on $\mathcal{B}$.
\end{proof}


	
\subsection{Variety and its Dimension}
One of the important notions we will use in our proofs is that of an affine algebraic variety. Given a set of polynomials, the algebraic variety defined by these polynomials is defined to be the set of their common zeros. That is, if $S$ is a set of polynomials on $n$ variables over a field $\F$, then
\[\mathcal V=\{\mathbf a\in\F^n:\forall f\in S,f(\mathbf a)=0\}.\]
Given a variety, an important property that is studied is its dimension. Intuitively, it is an appropriate generalisation of the notion of dimension for linear spaces. We will not be defining it formally here and refer the interested reader to the book by Cox, Little and O'Shea \cite{CLO}. However we state formally the properties, of dimensions of varieties, that we will be using in our proofs.
\begin{lemma}[Section 2.8 in {\cite{smith2014introduction}}]\label{lem:dimlower}
Let $S$ be a set of polynomials in $n$ variables over an algebraically closed field $\F$ such that $|S|\le n$. Let $\mathcal V=\mathbb V(S)$ be the set of common zeros of polynomials in $S$. If $\mathcal V$ is non-empty, then the dimension of $\mathcal V$ is at least $n-|S|$.
\end{lemma}
\begin{lemma}[Chapter 4 in {\cite{CLO}}]\label{lem:diminclusion}
Let $\F$ be an algebraically closed field, and let $\mathcal V_1\subseteq\F^n$ and $\mathcal V_2\subseteq\F^n$ be two affine varieties such that $\mathcal V_1\subseteq\mathcal V_2$. Then, the dimension of $\mathcal V_1$ is at most the dimension of $\mathcal V_2$.
\end{lemma}

\section{A lower bound for Algebraic Branching Programs}	\label{sec:layered}
In this section we prove \autoref{thm:intro main}. We start by restating it.

\begin{theorem}\label{thm:main}
	Let $n \in \N$ and let $\F$ be a field such that $\Char(\F) \nmid n$.
	If $\mathcal{A}$ is an algebraic branching program with edge labels of degree at most $\Delta$ that computes the polynomial $\sum_{i=1}^{n} x_i^n$, then the size of $\mathcal{A}$ is at least
	\[
		\Omega \inparen{\frac{n^2}{\Delta}}.
	\]
\end{theorem}
	
For technical reasons, we work with a slightly more general model which we call \emph{multilayered} ABPs, which we now define. 
	
\begin{definition}[Multilayered ABP]
	Let $\mathcal{A}_1, \ldots, \mathcal{A}_k$ be $k$ ABPs with $d_1, \ldots, d_k$ layers and $\tau_1, \ldots, \tau_k$ vertices, respectively. A multilayered ABP $\mathcal{A}$, denoted by $\mathcal{A} = \sum_{i = 1}^k \mathcal{A}_i$, is the ABP obtained by placing $\mathcal{A}_1, \mathcal{A}_2, \ldots, \mathcal{A}_k$ in parallel and identifying their start and end vertices respectively. Thus, the polynomial computed by $\mathcal{A}$ is $\sum_{i=1}^k [\mathcal{A}_i]$, where $ [\mathcal{A}_i]$ is the polynomial computed by $\mathcal{A}_i$. 
	
The number of layers of $\mathcal{A}$ is $d := \max \set{d_1, \ldots, d_k}$. The size of $\mathcal{A}$ is the number of vertices in $\mathcal{A}$, and thus equals  \[\abs{ \mathcal{A}} := 2 + \sum_i (\tau_i-2) \, .\] 
\end{definition}

This model is an intermediate model between (layered) ABPs and unlayered ABPs: given a multilayered ABP of size $\tau$ it is straightforward to construct an unlayered ABP of size $O(\tau)$ which computes the same polynomial. 	

	
\subsection{A robust lower bound for ABPs of formal degree at most $n$}
In this section, we prove a lower bound for the case where the formal degree of every vertex in the ABP is at most $n$. 
In fact, Kumar \cite{K19} has already proved a quadratic lower bound for this case.
	
\begin{theorem}[\cite{K19}]\label{thm:homog}
	Let $n \in \N$ and let $\F$ be a field such that $\Char(\F) \nmid n$. Then any algebraic branching program of formal degree at most $n$ which computes the polynomial $\sum_{i = 1}^n x_i^n$ has at least $\Omega(n^2)$ vertices. 
\end{theorem}
	
However, to prove \autoref{thm:main}, we need the following more ``robust'' version of \autoref{thm:homog}, which gives a lower bound for a larger class of polynomials. For completeness, we also sketch an argument for the proof which is a minor variation of the proof of \autoref{thm:homog}.

\begin{theorem}\label{thm:robustHomog}
	Let $n \in \N$ and let $\F$ be field such that $\Char(\F) \nmid n$.
	Let $A_1(\vx), \ldots, A_r(\vx), B_1(\vx),  \ldots, B_r(\vx)$ and $R(\vx)$ be polynomials such that for every $i$, $A_i(\mathbf{0}) = B_i(\mathbf{0}) = 0$ and $R$ is a polynomial of degree at most $n-1$. Then, any algebraic branching program over $\F$, of formal degree at most $n$ and edge labels of degree at most $\Delta \le n/10$, which computes the polynomial 
	\[\sum_{i = 1}^n x_i^n + \sum_{j = 1}^r A_j\cdot B_j + R\]
	has at least $\frac{(n/2-r)n}{2\Delta}$ vertices. 
\end{theorem}
	

The proof of the theorem follows from \autoref{lem:structure} and the following lemma which is a slight generalization of  Lemma 3.1 in \cite{K19}. We include a proof for completeness. 


\sloppy
\begin{lemma}\label{lem:varietyDimLB}
	Let $n \in \N$, and let $\F$ be an algebraically closed field such that $\Char(\F) \nmid n$. Let $\set{P_1, \ldots, P_m, Q_1, \ldots, Q_m, A_1, \ldots, A_r, B_1, \ldots, B_r}$ be a set of polynomials in $\F[x_1, \ldots, x_n]$ such that the set of their common zeros 
	\[
		V = \mathbb{V}(P_1, \ldots, P_m, Q_1, \ldots, Q_m, A_1, \ldots, A_r, B_1, \ldots, B_r) \subseteq \F^n
	\]
	is non-empty. Finally, suppose $R$ is a polynomial in $\F[\vx]$ of degree at most $n-1$, such that
	\[
		\sum_{i=1}^{n} x_i^n + \sum_{j=1}^{r} A_j \cdot B_j = R + \sum_{i=1}^{m} P_i \cdot Q_i.
	\]
	Then, $m \geq \frac{n}{2} - r$.
\end{lemma}
		
\begin{proof}
	Since $V\neq\emptyset$, by \autoref{lem:dimlower}, $\dim(V)\ge n-2m-2r$.  Thus, the set of zeros with multiplicity two of 
	\[
		\sum_{i=1}^{n} x_i^n - R = \sum_{i=1}^{m} P_i \cdot Q_i - \sum_{j=1}^{r} A_j \cdot B_j, 
	\]
	has dimension at least $n - 2m - 2r$. Now if $S$ is the set of common zeros of the set of all first order partial derivatives of $\sum_{i=1}^n x_i^n-R$, $V\subseteq S$. Thus, by \autoref{lem:diminclusion}, $\dim(S)\ge n-2m-2r$. Up to scaling by $n$ (which is non-zero in $\F$, by assumption), the set of all first order partial derivatives of $\sum_{i=1}^{n} x_i^n - R$ is given by 
	\[
		\set{x_i^{n-1} - \frac{1}{n}\partial_{x_i}R}_{i \in [n]}.
	\]
	Thus, the statement of this lemma immediately follows from the following claim.
			
	\begin{claim}[Lemma 3.2 in \cite{K19}]
		Let $\F$ be an algebraically closed field, and $D$ a positive natural number. For every choice of polynomials $g_1, g_2, \ldots g_n \in \F[\vx]$ of degree at most $D-1$, the dimension of the variety 
		\[\mathbb{V}(x_1^D - g_1, x_2^D - g_2, \ldots, x_n^D - g_n)\] is zero.
	\end{claim}
	Indeed, the above claim shows that $0 = \dim(S) \geq n-2m-2r$, and so $m \geq \frac{n}{2} - r$. This completes the proof of \autoref{lem:varietyDimLB}.
\end{proof}

We now use \autoref{lem:structure} and \autoref{lem:varietyDimLB} to complete the proof of \autoref{thm:robustHomog}.

\begin{proof}[Proof of \autoref{thm:robustHomog}]
	Let $\mathcal{B}$ be an algebraic branching program of formal degree at most $n$, edge labels of degree at most $\Delta \le n/10$, and with start vertex $s$ and end vertex $t$, which computes
	\[\sum_{i = 1}^n x_i^n + \sum_{j = 1}^r A_j\cdot B_j + R.\]
	We may assume without loss of generality that $\F$ is algebraically closed, by interpreting $\mathcal{B}$ as an ABP over the algebraic closure of $\F$, if necessary.
		
	Let $n' = \lfloor n/\Delta \rfloor$, fix $k \in \set{1, 2, \ldots ,n'-1}$, and let $V_k = \set{v_{k,1}, v_{k,2}, \ldots, v_{k,m}}$ be the set of all vertices in $\mathcal{B}$ whose formal degree lies in the interval $[k\Delta, (k+1)\Delta)$.  Letting $P'_j = [s,v_{k,j}]$, by \autoref{lem:structure}, there exist polynomials $Q'_1, Q'_2, \ldots, Q'_m$ and $R'$, each of degree at most $n-1$ such that
	\[
		\sum_{i = 1}^n x_i^n + \sum_{j = 1}^r A_j\cdot B_j + R = \sum_{j=1}^{m} P'_j \cdot Q'_j + R'\, .
	\]
	Let $\alpha_j$, $\beta_j$ be the constant terms in $P'_j$, $Q'_j$ respectively. Then by defining
	\[P_j = P'_j - \alpha_j \qquad \text{and} \qquad Q_j = Q'_j - \beta_j,\]
	we have that
	\[
		\sum_{i=1}^{n} x_i^n + \sum_{j=1}^{r} A_j \cdot B_j = R'' + \sum_{j=1}^{m} P_j \cdot Q_j.
	\]
	Here, $R'' = -R + R' + \sum_{j=1}^{m} (\alpha_j \cdot Q'_j + \beta_j \cdot P'_j + \alpha_j \beta_j)$. We now have that for every $i$, the constant terms of $P_i$, $Q_i$ are zero and $\deg(R'') \leq n-1$.
	Let 
	\[
		\mathcal{V} = \mathbb{V}(P_1, \ldots, P_m, Q_1, \ldots, Q_m, A_1, \ldots, A_r, B_1, \ldots, B_r).
	\]
	Then $\mathbf{0} \in \mathcal{V}$, and so $\mathcal{V} \neq \emptyset$. Thus by \autoref{lem:varietyDimLB}, we know that $m \geq \frac{n}{2} - r$.
		
	Finally, for $ k\neq k' \in \set{1, 2, \ldots, n'-1}$, $V_k \cap V_{k'} = \emptyset$. Thus, the number of vertices in $\mathcal{B}$ must be at least
	\[
		\inparen{\dfrac{n}{2} - r} \cdot (n'-1) \geq \inparen{\frac{n}{2} - r}\cdot  \frac{n}{2\Delta}\,.
		\qedhere
	\]
\end{proof}
	

\subsection{A lower bound for the general case}

The following lemma shows how we can obtain, given an ABP with $d$ layers which computes a polynomial $F$, a multilayered ABP, whose number of layers is significantly smaller, which computes $F$ plus a small ``error term''.
	
\begin{lemma}\label{lem:shrinkage for a single layered abp}
	Let $\mathcal{A}$ be an ABP over a field $\F$ with $d$ layers, which computes the polynomial $F$ and has $m$ vertices. Let $s$ and $t$ be the start and end vertices of $\mathcal{A}$ respectively, and let $L = \set{u_1, u_2, \ldots, u_{|L|}}$ be the set of vertices in the $\ell$-th layer of $\mathcal{A}$. For every $i \in \set{1, 2, \ldots, |L|}$, let $\alpha_i$ and $\beta_i$ be the constant terms of $[s,u_i]$ and $[u_i,t]$ respectively. Furthermore, let $P_i$ and $Q_i$ be polynomials such that $[s,u_i] = P_i + \alpha_i$ and $[u_i,t] = Q_i + \beta_i$.
		
	Then, there is a multilayered ABP ${\mathcal{{A}}}'$,  with at most $\max\{\ell, d-\ell + 1\}$ layers and size at most $\abs{\mathcal{A}}$ that computes the polynomial 
	\[
		F - \sum_{i=1}^{\abs{L}} P_i \cdot Q_i + \sum_{i=1}^{\abs{L}} \alpha_i \cdot \beta_i .
	\]
\end{lemma}
	
\begin{proof}
	Let $u_1, u_2, \ldots, u_{|L|}$ be the vertices in $L$ as described, so that
	\[
		F = [s,t] = \sum_{i = 1}^{|L|} [s, u_{i}]\cdot [u_{i}, t] \, . 
	\]
	Further, for every $i \in \set{1, 2, \ldots, |L|}$, $[s,u_i] = P_i + \alpha_i$ and $[u_i,t] = Q_i + \beta_i$, where the constant terms of $P_i$ and $Q_i$ are zero (by definition). Having set up this notation, we can thus express the polynomial $F$ computed by $\mathcal{A}$ as 
	\[
		F = [s,t] = \sum_{i = 1}^{|L|} (P_i + \alpha_i)\cdot (Q_i + \beta_i) \, .
	\]
	On further rearrangement, this gives 
	\[
		F - \inparen{\sum_{i = 1}^{|L|} P_i \cdot Q_i} + \inparen{\sum_{i = 1}^{|L|} \alpha_i\cdot \beta_i} = \inparen{\sum_{i = 1}^{|L|} \alpha_i \cdot (Q_i + \beta_i)} + \inparen{\sum_{i = 1}^{|L|} (P_i + \alpha_i)\cdot \beta_i} \, .
	\] 
	This is equivalent to the following expression.
	\[
		F - \inparen{\sum_{i = 1}^{|L|} P_i \cdot Q_i} + \inparen{\sum_{i = 1}^{|L|} \alpha_i\cdot \beta_i} =  \inparen{\sum_{i = 1}^{|L|} \alpha_i\cdot [u_i, t]} + \inparen{\sum_{i = 1}^{|L|} [s, u_i]\cdot \beta_i} \, .
	\] 
	Now, observe that the polynomial $\sum_{i=1}^{|L|} [s, u_i] \cdot \beta_i$ is computable by an ABP $\mathcal{B}$ with $\ell+1$ layers, obtained by just keeping the vertices and edges within first $\ell$ layers of $\mathcal{A}$ and the end vertex $t$, deleting all other vertices and edges, and connecting the vertex $u_i$ in the $\ell$-th layer to $t$ by an edge of weight $\beta_i$.  
	Similarly, the polynomial $\sum_{i=1}^{|L|} \alpha_i\cdot [u_i, t]$ is computable by an ABP $\mathcal{C}$ with at most $(d-\ell + 1)+1$ layers,  whose set of vertices is $s$ along the vertices in the layers $\ell, \ell + 1, \ell+2, \ldots, d$ of $\mathcal {A}$. From the definition of $\mathcal{B}$ and $\mathcal{C}$, it follows that the multilayered ABP $\tilde{\mathcal{A}}$ obtained by taking the sum of $\mathcal{B}$ and $\mathcal{C}$ has at most $\max\set{\ell+1, d-\ell+2}$ layers. 
		
	We are almost done with the proof of the lemma, except for the upper bound on the number of vertices of the resulting multilayered ABP $\tilde{\mathcal{A}}$, and the fact that the upper bound on the depth is slightly weaker than claimed. Both these issues can be solved simultaneously.
	
	 The vertices in $L$ appear in both the ABP $\mathcal{B}$ and the ABP $\mathcal{C}$ and are counted twice in the size of $\tilde{\mathcal{A}}$. However, every other vertex is counted exactly once.  Hence, 
	\begin{equation}\label{eqn:overlap}
		\abs{\mathcal{B}} + \abs{\mathcal{C}} =  \abs{\mathcal{A}} + \abs{L} \, .
	\end{equation}
 	In order to fix this issue, we first observe that the edges between the vertices in the $\ell$-th layer of $\mathcal{B}$ and the end vertex $t$ are labeled by $\beta_{1}, \beta_2, \ldots, \beta_{|L|}$, all of which are field constants. In the following claim, we argue that for ABPs with this additional structure, the last layer is redundant and can be removed.

	\begin{claim}\label{clm:last layer removal}
		Let ${\mathcal{M}}$ be an ABP over $\F$ with $k+1$ layers and edge labels of degree at most $\Delta$ such that the labels of all the edges between the $k$-th layer of ${\cal M}$ and its end vertex are scalars in $\F$. Then, there is an ABP $\mathcal{M}'$ with  $k$ layers computing the same polynomial as $\mathcal{M}$, with edge labels of degree at most $\Delta$, such that
		\[\abs{\mathcal{M}'} \leq \abs{\mathcal{M}} - \abs{V} \, , \]
		where $V$ is the set of vertices in the $k$-th layer of $\mathcal{M}$.
	\end{claim}
	
	An analogous statement, with an identical proof, is true if we assume that all edge labels between the first and second layer are scalars in $\F$.

	We first use \autoref{clm:last layer removal} to complete the proof of the lemma. As observed above, the edge labels between the last layer $L$ of $\mathcal{B}$ and its end vertex are all constants. Hence, by \autoref{clm:last layer removal}, there is an ABP $\mathcal{B}'$ which computes the same polynomial as $\mathcal{B}$ such that $\abs{\mathcal{B}'} \leq \abs{\mathcal{B}} - \abs{L}$, and $\mathcal{B}'$ has only $\ell$ layers. Similarly, we can obtain an ABP $\mathcal{C}'$ with at most $d-\ell+1$ layers.
	
	We consider the multilayered ABP $\mathcal{A}'$ by taking the sum of $\mathcal{B}'$ and $\mathcal{C}'$. Clearly, the number of layers in $\mathcal{A}'$ is at most $\max\{\ell, d-\ell+1\}$ and the size is at most 
	\[\abs{\mathcal{A}'} \leq  \abs{\mathcal{B}'} + \abs{\mathcal{C'}} \leq  \inparen{\abs{\mathcal{B}} - \abs{L}} + \inparen{\abs{\mathcal{C} }- \abs{L}} \leq \abs{\mathcal{A}} \, .\]
	Here, the second inequality follows by \autoref{clm:last layer removal} and the last one follows by \autoref{eqn:overlap}. To complete the proof of the lemma, we now prove \autoref{clm:last layer removal}. 
\end{proof}
	
\begin{proof}[Proof of \autoref{clm:last layer removal}]
	For the proof of the claim, we  focus on the $k$-th and $(k-1)$-st layer of $\mathcal{M}$. To this end, we first set up some notation. Let $\set{v_1, v_2, \ldots, v_r}$ be the set of vertices in the $k$-th layer of $\mathcal{M}$, $\set{u_1, u_2, \ldots, u_{r'}}$ be the set of vertices in $(k-1)$-st layer of $\mathcal{M}$, and $a$, $b$ denote the start and the end vertices of $\mathcal{M}$ respectively. Then, the polynomial computed by $\mathcal{M}$, can be decomposed as 
	\[
		[a,b] = \sum_{i = 1}^r [a, v_i] \cdot  [v_i, b]\, .
	\]
	Note that $(v_i,b)$ is an edge in the ABP. Similarly, the polynomial $[a,v_i]$ can be written as 
	\[
		[a,v_i] = \sum_{j = 1}^{r'} [a, u_j] \cdot [u_j, v_i] \, .
	\]
	Combining the two expressions together, we get  
	\[
		[a,b] = \sum_{i = 1}^r [v_i, b] \cdot  \inparen{\sum_{j = 1}^{r'} [u_j,v_i] \cdot [a, u_j]}\, ,
	\]
	which on further rearrangement,  gives us 
	\begin{equation} \label{eqn:decomposition using second last layer}
		[a,b] = \sum_{j = 1}^{r'} \inparen{\sum_{i = 1}^{r} [v_i,b] [u_j, v_i]} \cdot [a, u_j] \, .
	\end{equation}
 	From the hypothesis of the claim, we know that for every $i \in [r]$, the edge label $[v_i,b]$ is a field constant, and the edge label $[u_j, v_i]$ is a polynomial of degree at most $\Delta$. Thus, for every $j \in [r']$, the expression $\inparen{\sum_{i = 1}^{r} [u_i,b] [u_j,v_i]}$ is a polynomial of degree at most $\Delta$.
 	
 	This gives us the following natural construction for the ABP $\mathcal{M}'$ from $\mathcal{M}$. We delete the vertices $v_1, v_2, \ldots, v_r$ in $\mathcal{M}$ (and hence, all edges incident to them), and for every $j \in \set{1, 2, \ldots, r'}$, we connect the vertex $u_j$ with the end vertex $b$ using an edge with label $\inparen{\sum_{i = 1}^{r} [v_i,b] [u_j, v_i]}$. The upper bound on the size and the number of layers of $\mathcal{M}'$ is immediate from the construction, and that it computes the same polynomial as $\mathcal{M}$ follows from \autoref{eqn:decomposition using second last layer}.
 \end{proof}	

We now state and prove a simple generalization of \autoref{lem:shrinkage for a single layered abp} for a multilayered ABP. 
	
\begin{lemma}\label{lem:shrinkage general}
	Let $\mathcal{A} = \sum_{i=1}^{m} \mathcal{A}_i$ be a multilayered ABP with $d$ layers over a field $\F$ computing the polynomial $F$, such that each $\mathcal{A}_i$ is an ABP with $d_i$ layers.
	Also, let $\ell_{i,j}$ be the number of vertices in the $j$-th layer of $\mathcal{A}_i$ ($\ell_{i,j} = 0$ if $\mathcal{A}_i$ has fewer than $j$ layers), and $\ell = \min_{j \in (d/3, 2d/3)} \set{\sum_{i=1}^m \ell_{i,j}}$.
	
	Then, there is a multilayered ABP  with at most $2d/3$ layers and size at most $\abs{\mathcal{A}}$ that computes a polynomial of the form
	\[F - \sum_{i=1}^{\ell} P_i \cdot Q_i + \delta \, , \]
	where $\set{P_1, \ldots, P_{\ell}, Q_1, \ldots, Q_{\ell}}$ is a set of non-constant polynomials with constant term zero and $\delta \in \F$.
\end{lemma}
	
\begin{proof}
	Let $j_0 \in (d/3, 2d/3)$ be the natural number which minimizes the quantity $\sum_{i=1}^m \ell_{i,j}$, and let $S \subseteq [m]$ be the set of all indices $i$ such that $\mathcal{A}_i$ has at least $j_0$ layers. Let $\mathcal{A}' = \sum_{i \in S}\mathcal{A}_i$ and $\mathcal{A}'' = \sum_{i \notin S}\mathcal{A}_i$. Thus, 
	\[
		\mathcal{A} = \mathcal{A}' + \mathcal{A}''.
	\]		
	Here, $\mathcal{A}'' = \sum_{i \notin S}\mathcal{A}_i$ is a multilayered ABP with at most $2d/3$ layers. Moreover, $\abs{\mathcal{A}} = \abs{\mathcal{A}'} + \abs{\mathcal{A}''}$.
		
	The idea now is to apply \autoref{lem:shrinkage for a single layered abp} to every ABP in $\mathcal{A}'$. For every $i \in S$, we know that there exist some polynomials $P_{i,1}, \ldots, P_{i, \ell_{i, {j_0}}}, Q_{i,1}, \ldots, Q_{i, \ell_{i, {j_0}}}$ with constant terms zero and a constant $\delta_i$, such that  
	\[
		F_i - \sum_{r = 0}^{\ell_{i, {j_0}}} P_{i, r}Q_{i,r} + \delta_i
	\]
	can be computed by a multilayered ABP. Let us denote this multilayered ABP by $\mathcal{B}_i$. From \autoref{lem:shrinkage for a single layered abp}, we know that $\mathcal{B}_i$ has at most $\max\{j_0, d_i-j_0+1\} \leq 2d/3$ layers and size at most $\abs{\mathcal{A}_i}$. Taking a sum over all $i \in S$ and re-indexing the summands, we get that there exist polynomials $P_{1}, \ldots, P_{\ell}, Q_{1}, \ldots, Q_{\ell}$ with constant terms zero and a constant $\delta$ such that the polynomial 
	\[
		\sum_{i\in S} F_i - \sum_{r = 0}^{\ell} P_{r}Q_{r} + \delta
	\]
	is computable by a multilayered ABP $\mathcal{B} = \sum_{i \in S} \mathcal{B}_i$ with at most $2d/3$ layers and size at most $\sum_{i \in S} \abs{\mathcal{A}_i} \leq \abs{\mathcal{A}'}$. Finally, by combining the multilayered ABPs $\mathcal{B}$ and $\mathcal{A}''$, we get that the polynomial 
	\[
		F  - \sum_{r = 0}^{\ell} P_{r}Q_{r} + \delta
	\]
	is computable by a multilayered ABP with at most $2d/3$ layers and size at most $\abs{\mathcal{A}}$.
\end{proof}
	
We now use \autoref{lem:shrinkage general} to prove our main result, which we restate once more.
\paragraph{Theorem 7 (restated).} Let $n\in\N$ and let $\F$ be a field such that $\Char(\F)\nmid n$. If $\mathcal A$ is an algebraic branching program with edge labels of degree at most $\Delta$ that computes the polynomial $\sum_{i=1}^n x_i^n$, then the size of $\mathcal A$ is at least $\Omega(n^2/\Delta)$.
 
	
\begin{proof}[Proof of \autoref{thm:main}]
	Let $\mathcal{A}$ be a multilayered ABP with $d_0$ layers which computes the polynomial $\sum_{i = 1}^n x_i^n$. As before we may assume without loss of generality that the underlying field $\F$ is algebraically closed. Note that if $d_0$ is at most $n/\Delta$, then the formal degree of $\mathcal A$ is at most $d_0\Delta\le n$. Thus, by \autoref{thm:robustHomog}, we know that $\abs{\mathcal A}$ is at least $\Omega(n^2/\Delta)$ and we are done. Also, if $d_0 > n^2/\Delta$, then again we have our lower bound since each layer of $\mathcal{A}$ must have at least one vertex. Thus, we can assume that $n/\Delta \le d_0 \le n^2/\Delta$. 
		
	The proof idea is to iteratively make changes to $\mathcal{A}$ till we get a multilayered ABP $\mathcal{A}'$ of formal degree at most $n$ that computes a polynomial of the type 
	\[\sum_{i = 1}^n x_i^n + \sum_{j = 1}^r A_j\cdot B_j + R\]
	where $r \leq n/10$ and $A_1(\vx), \ldots, A_r(\vx), B_1(\vx), \ldots, B_r(\vx), R(\vx)$ are polynomials such that for every $i$, $A_i(\mathbf{0}) = B_i(\mathbf{0}) = 0$ and $R$ has degree at most $n-1$. 
	Once we have this, we can invoke \autoref{thm:robustHomog} and get the required lower bound.
		
	We now explain how to iteratively obtain $\mathcal{A}'$ from $\mathcal{A}$. In one step, we ensure the following.
		
	\begin{claim}\label{clm:oneStep}
		Let $\mathcal{A}_k$ be a multilayered ABP with edge labels of degree at most $\Delta$, $d_k \ge n/\Delta$ layers and size at most $\tau$ that computes a polynomial of the form $\sum_{i = 1}^n x_i^n + \sum_{j = 1}^r A_j\cdot B_j + R$ where $A_1(\vx), \ldots, A_r(\vx), B_1(\vx), \ldots, B_r(\vx), R(\vx)$ are polynomials such that for every $j$, $A_j(\mathbf{0}) = B_j(\mathbf{0}) = 0$ and $R$ has degree at most $n-1$.
			
		If $\tau \leq 0.001 n^2/\Delta$, then there exists a multilayered ABP $\mathcal{A}_{k+1}$ with at most $2 d_k/3$ layers and size at most $\tau$ which computes a polynomial of the form 
		\[\sum_{i = 1}^n x_i^n + \sum_{j = 1}^{r'} A'_j\cdot B'_j + R'\, ,\]
		such that $r' \leq  r + 0.005\frac{n^2}{\Delta \cdot d_k}$ and $A'_1(\vx), \ldots, A'_{r'}(\vx), B'_1(\vx), \ldots, B'_{r'}(\vx), R'(\vx)$ are polynomials such that for every $i$, $A'_i(\mathbf{0}) = B'_i(\mathbf{0}) = 0$ and $R'$ has degree at most $n-1$.
	\end{claim}
	
	Before moving on to the proof of  \autoref{clm:oneStep}, we first use it to complete the proof of  \autoref{thm:main}. Let us set $\mathcal{A}_0 = \mathcal{A}$. Then, $\mathcal{A}_0$ is a multilayered ABP with $d_0$ layers and size at most $\tau $ that computes the polynomial $\sum_{i = 1}^n x_i^n$.
		
	If $\tau \geq 0.001 n^2/\Delta$, the statement of the theorem follows. Otherwise, we apply \autoref{clm:oneStep} iteratively $K$ times, as long as the number of layers is more than $n/\Delta$, to  eventually get a multilayered ABP $\mathcal{A}' = \mathcal{A}_{K}$ with $d' \le n/\Delta$ layers. Let $d_0, \ldots, d_{K-1},d_K$ denote the number of layers in each ABP in this sequence, so that $d_{K-1} > n/\Delta$, and $d_k \le 2 d_{k-1}/3$ for $k \in [K]$.  $\mathcal{A}'$ is an ABP with at most $n/\Delta$ layers and size at most $\tau$, which by induction, computes a polynomial of the form
	\[\sum_{i = 1}^n x_i^n + \sum_{j = 1}^r A_j\cdot B_j + R \, ,\]
	where $A_1(\vx), \ldots, A_r(\vx), B_1(\vx), \ldots, B_r(\vx), R(\vx)$ are polynomials such that for every $i$, $A_i(\mathbf{0}) = B_i(\mathbf{0}) = 0$ and $R$ has degree at most $n-1$. Further, the number of error terms, $r$, is at most 
	\[
 		\frac{0.005 n^2} {\Delta} \left( \frac{1}{d_{K-1}} + \frac{1}{d_{K-2}} + \cdots + \frac{1}{d_0} \right).
	\]
	Since $d_k \le \frac{2}{3} \cdot d_{k-1}$, we have that $\frac{1}{d_{k-1}} \le \frac{2}{3} \cdot \frac{1}{d_k}$ for all $k \in [K]$, so that
\[
r \le \frac{0.005n^2}{\Delta} \cdot \frac{1}{1-2/3} \cdot \frac{1}{d_{K-1}} \le \frac{n}{10}
\]
as $d_{K-1} \ge n/\Delta$.	

	At this point, since the formal degree is at most $n$, using \autoref{thm:robustHomog} we get 
	\[
		\tau \geq \abs{\mathcal{A}'} \geq \dfrac{(n/2-r)n}{2\Delta} = \Omega \inparen{\dfrac{n^2}{\Delta}} \, . \qedhere 
	\]
\end{proof}
		
To complete the proof of  \autoref{thm:main}, we now prove \autoref{clm:oneStep}.

\begin{proof}[Proof of \autoref{clm:oneStep}]
	Let $\mathcal{A}_k = \sum_{i=1}^{m} \mathcal{A}_{k,i}$, and for $j \in [d_k]$, let $\ell_{i,j}$ be the number of vertices in layer $j$ of $\mathcal{A}_{k,i}$. Recall that if the number of layers in $\mathcal{A}_{k,i}$ is strictly less than $j$, then we set $\ell_{i,j} = 0$. Let $\ell$ be the total number of vertices in the middle layers of $\mathcal{A}_k$, defined as 
	\[ \ell = \sum_{i=1}^{m} \inparen{\sum_{j \in (d_k/3, 2d_k/3)} \ell_{i,j}} \, .\] 
	Since $\ell \leq \tau \leq \frac{0.001 n^2}{\Delta }$, by averaging, we know that there is a $j_0 \in (d_k/3, 2d_k/3)$, such that 
	\[
		\ell_{j_0} = \sum_{i=1}^{m}\ell_{i,{j_0}} \leq \frac{\ell}{d_k/3} \leq \frac{0.001 n^2}{\Delta } \cdot \frac{1}{d_k/3} \leq 0.005\frac{n^2}{\Delta \cdot d_k} \, .
	\]
	This condition, together with \autoref{lem:shrinkage general}, tells us that there is a multilayered ABP  $\mathcal{A}_{k+1}'$ with at most $2d_k/3$ layers and size at most $\tau$ that computes a polynomial of the form
	\[
		\sum_{i = 1}^n x_i^n + \sum_{j = 1}^r A_j\cdot B_j + R - \sum_{i=1}^{\ell_{j_0}} P_i \cdot Q_i + \delta \, , 
	\]
	where  $P_1, \ldots, P_{\ell}, Q_1, \ldots, Q_{\ell}$ are a set of non-constant polynomials with constant term zero and $\delta \in \F$. Since $\ell_{j_0} \leq 0.005\frac{n^2}{\Delta \cdot d_k}$, the claim follows.	
\end{proof}



\begin{thebibliography}{22}
\bibitem[6]{CLO} David A. Cox, John B. Little, and Donal O’Shea. Ideals, Varieties and Algorithms. Under- graduate texts in mathematics. Springer, 2007. doi:10.1007/978-0-387-35651-8.
\bibitem[13]{K19} Mrinal Kumar. A quadratic lower bound for homogeneous algebraic branching programs. Computational Complexity, 28(3):409–435, 2019. doi:10.1007/s00037-019-00186-3.
\bibitem[22]{smith2014introduction} Justin R. Smith. Introduction to Algebraic Geometry. Textbooks in Mathematics. Taylor \& Francis, 2014. URL: https://books.google.com/books?id=zx7usgEACAAJ.
\end{thebibliography}
\end{document}
